\section{Methodology}
\label{sec:methodology}

\subsection{Data Construction}

The numerical data are generated according to the instructions provided
in the assignment. The student's enrollment number is converted into
eight integer digits, replacing every zero with one. These digits are
then used to construct the coefficient matrix $A$, the vector
$\mathbf{b}$, and the experimental observations $(x_i,y_i)$.

The original data-generation expressions are preserved. NumPy is used
for the array construction specified in the assignment, whereas the
principal linear algebra operations are implemented explicitly using
Python loops. In particular, no predefined routines for QR
factorization, least-squares solutions, or polynomial fitting are used.

\subsection{QR Factorization by Modified Gram--Schmidt}

Let $B\in\mathbb{R}^{m\times n}$, with $m\geq n$, be a matrix with
linearly independent columns. The QR factorization expresses $B$ as
\begin{equation}
    B=QR,
    \label{eq:qr-factorization}
\end{equation}
where $Q\in\mathbb{R}^{m\times n}$ has orthonormal columns,
$Q^TQ=I_n$, and $R\in\mathbb{R}^{n\times n}$ is upper triangular.

The factorization is computed using the modified Gram--Schmidt
(MGS) algorithm. This method constructs an orthonormal basis for the
column space of $B$ by successively removing the components of each
column in the directions of previously computed basis vectors.

Denote the $j$-th column of $B$ by $\mathbf{b}_j$. Initialize
\begin{equation}
    \mathbf{v}_j^{(0)}=\mathbf{b}_j.
\end{equation}
For $i=1,\ldots,j-1$, compute
\begin{align}
    r_{ij}
     & =\mathbf{q}_i^T\mathbf{v}_j^{(i-1)},      \\
    \mathbf{v}_j^{(i)}
     & =\mathbf{v}_j^{(i-1)}-r_{ij}\mathbf{q}_i.
\end{align}
The next diagonal entry of $R$ and the corresponding column of $Q$
are then obtained as
\begin{align}
    r_{jj} & =\|\mathbf{v}_j^{(j-1)}\|_2,          \\
    \mathbf{q}_j
           & =\frac{\mathbf{v}_j^{(j-1)}}{r_{jj}}.
\end{align}

The algorithm is applied first to the matrix $A$ and subsequently to
the design matrix $X$. The resulting factors are used to solve both
least-squares problems. During implementation, a diagonal entry of
$R$ that is numerically negligible is treated as an indication of
possible rank deficiency; the corresponding case must be reported
rather than dividing by that entry.

\subsection{Selection of the Orthogonalization Method}

Several methods can compute the QR factorization. The choice depends
on numerical stability, implementation complexity, and the objectives
of the computational assignment.

\begin{table}[htbp]
    \centering
    \small
    \caption{Comparison of methods for computing the QR factorization.}
    \label{tab:qr-methods}
    \begin{tabular}{p{2.5cm} p{5.3cm} p{5.3cm}}
        \hline
        \textbf{Method} & \textbf{Advantages} & \textbf{Limitations}       \\
        \hline
        Classical Gram--Schmidt
                        &
        Simple formulation and straightforward implementation.
                        &
        Can lose orthogonality when the columns are nearly linearly
        dependent.                                                         \\[0.3em]

        Modified Gram--Schmidt
                        &
        Explicit construction of the orthonormal basis; generally preserves
        orthogonality better than the classical variant and is convenient
        for implementing the algorithm from first principles.
                        &
        Can still lose orthogonality for ill-conditioned matrices; it is
        not generally as numerically robust as Householder QR.             \\[0.3em]

        Householder reflections
                        &
        Highly reliable numerical behavior and a standard choice for dense
        QR factorizations.
                        &
        Requires implementing and managing reflection transformations,
        making the algorithm less direct to explain from first principles. \\[0.3em]

        Givens rotations
                        &
        Useful for exploiting sparsity and updating factorizations when
        individual rows or entries change.
                        &
        Requires managing a sequence of rotations, which is unnecessary
        for the small dense matrices considered here.                      \\
        \hline
    \end{tabular}
\end{table}

The modified Gram--Schmidt method is selected primarily for its
transparent computational structure and suitability for a
from-scratch implementation. It provides a direct illustration of
orthogonalization while improving numerical behavior relative to
classical Gram--Schmidt. Householder QR would generally be preferable
when numerical robustness is the primary concern, particularly for
larger or ill-conditioned problems
\cite{golub2013matrix,trefethen1997numerical,higham2002accuracy}.

The choice is therefore pedagogical and computational rather than
a claim that modified Gram--Schmidt is universally the most stable
QR algorithm.

\subsection{Least-Squares Solution Using QR}

Consider the overdetermined system
\begin{equation}
    B\mathbf{z}\approx\mathbf{c},
\end{equation}
where $B$ has full column rank. The least-squares solution minimizes
the squared Euclidean residual:
\begin{equation}
    \mathbf{z}^{*}
    =\underset{\mathbf{z}\in\mathbb{R}^{n}}{\arg\min}
    \,\|B\mathbf{z}-\mathbf{c}\|_2^2.
\end{equation}

Substituting $B=QR$ and using the orthonormality of the columns of
$Q$ gives
\begin{equation}
    \begin{aligned}
        \|B\mathbf{z}-\mathbf{c}\|_2^2
         & =
        \|R\mathbf{z}-Q^T\mathbf{c}\|_2^2 \\
         & \quad+
        \|(I-QQ^T)\mathbf{c}\|_2^2.
    \end{aligned}
    \label{eq:ls-orthogonal-decomposition}
\end{equation}

The second term is independent of $\mathbf{z}$. Consequently, the
least-squares solution is obtained by solving the upper triangular
system
\begin{equation}
    R\mathbf{z}=Q^T\mathbf{c}.
    \label{eq:triangular-system}
\end{equation}

This system is solved by back substitution, starting from the last
equation:
\begin{equation}
    z_i=
    \frac{
        d_i-\displaystyle\sum_{j=i+1}^{n}R_{ij}z_j
    }{R_{ii}},
    \qquad i=n,n-1,\ldots,1,
    \label{eq:back-substitution}
\end{equation}
where $\mathbf{d}=Q^T\mathbf{c}$.

For the first problem, $B=A$, $\mathbf{c}=\mathbf{b}$, and
$\mathbf{z}=\mathbf{x}$. For the second problem, $B=X$,
$\mathbf{c}=\mathbf{y}$, and $\mathbf{z}=[a,b,c]^T$.

The QR-based approach is preferred to solving the normal equations
$B^TB\mathbf{z}=B^T\mathbf{c}$. For a full-column-rank matrix, the
normal equations satisfy
$\kappa_2(B^TB)=\kappa_2(B)^2$, potentially amplifying the effects
of poor conditioning. The QR formulation avoids explicitly forming
$B^TB$ and reduces the least-squares problem to a triangular solve
\cite{golub2013matrix,higham2002accuracy}.

\subsection{Numerical Verification and Error Evaluation}

The computed QR factors are assessed using two quantities. The
relative factorization error is defined as
\begin{equation}
    \eta_{\mathrm{QR}}
    =
    \frac{\|B-QR\|_F}{\|B\|_F},
    \label{eq:qr-error}
\end{equation}
where $\|\cdot\|_F$ denotes the Frobenius norm. The orthogonality
error is defined as
\begin{equation}
    \eta_Q=\|Q^TQ-I_n\|_F.
    \label{eq:orthogonality-error}
\end{equation}
Both quantities are evaluated using explicitly implemented matrix
operations.

For the first problem, the residual norm
\begin{equation}
    \rho_A=\|A\mathbf{x}^{*}-\mathbf{b}\|_2
    \label{eq:linear-residual}
\end{equation}
is calculated to assess the least-squares solution. The uniqueness
of the solution is examined through the column rank of $A$ and the
diagonal entries of $R$.

For the quadratic fit, the predicted values are calculated from
\begin{equation}
    \widehat{y}_i=ax_i^2+bx_i+c,
\end{equation}
and the mean squared error is computed as
\begin{equation}
    \operatorname{MSE}
    =\frac{1}{6}\sum_{i=1}^{6}(y_i-\widehat{y}_i)^2.
\end{equation}

Finally, the experimental observations and the fitted polynomial
are plotted together. The numerical algorithms are implemented
without predefined linear algebra or polynomial-fitting routines;
a plotting library is used only to visualize the results.

The complete computational implementation, including the
data-generation code and the routines for QR factorization,
back substitution, error evaluation, and plotting, is provided
in the appendix.